% Part: proof-theory
% Chapter: natural-deduction
% Section: grafting

\documentclass[../../../include/open-logic-section]{subfiles}

\begin{document}

\olfileid{pt}{nat}{gra}

\olsection{پیوندزدن \usetoken{P}{proof}}

!!^{proof}ها در \Log{N1c} و~\Log{N1i}، !!{proof}هایی برای
!!{formula}های~$!A$ از فرض‌های~$!B$ هستند. فرض کنید خود~$!B$
!!a{proof} داشته باشد. چگونه !!{proof}~$!B$ را با !!{proof}~$!A$ از
$!B$ ترکیب کنیم تا !!a{proof} برای~$!A$ به دست آید که از فرض~$!B$
استفاده نکند؟

یک راه این است که با اعمال~\Intro{\lif} فرض~$!B$ را از !!{proof}~$!A$
حذف کنیم. سپس می‌توانیم !!{proof} حاصل برای~$!B \lif !A$ را با
!!{proof}~$!B$ ترکیب کنیم و به این برسیم:
\[
\AxiomC{$\Discharge{!B}{x}$}
\DeduceC{$!A$}
\DischargeRule{\Intro{\lif}}{x}
\UnaryInfC{$!B \lif !A$}
\AxiomC{}
\DeduceC{$!B$}
\RightLabel{\Elim{\lif}}
\BinaryInfC{$!A$}
\DisplayProof
\]
این کار سرراست، اما کمی نازیباست: چرا $\lif$ را معرفی کنیم تا بی‌درنگ
آن را حذف کنیم؟ باید بتوان از این نوع «میان‌بُرِ اضافی» از راه
$!B \lif !A$ پرهیز کرد. (اینکه همیشه می‌توان از چنین میان‌بُرهایی
پرهیز کرد، در واقع محتوای قضیهٔ نرمال‌سازی است.)

در \Log{N1c} و~\Log{N1i} در این مورد می‌توانیم به‌آسانی از آن پرهیز
کنیم: کافی است فرض‌های~$!B^x$ را با کل !!{proof}~$!B$ جایگزین کنیم.
این «پیوندزدن» !!{proof}ها به فرض‌های برهان‌های دیگر بسیار سودمند است؛
پس آن را در قالب یک تعریف ثبت می‌کنیم.

\begin{defn}
اگر $\delta_1$ یک \Log{N1c}-!!{proof} (یا \Log{N1i}-!!{proof}) برای
~$!A$ از فرض باز~$!B^x$ باشد، و $\delta_2$ یک
\Log{N1c}-!!{proof} (یا \Log{N1i}-!!{proof}) برای~$!B$ باشد، آنگاه
$\Subst{\delta_1}{\delta_2}{!B^x}$ حاصل جایگزینی هر رخدادِ مرخص‌نشدهٔ
~$!B^x$ در~$\delta_1$ با~$\delta_2$ است.
\end{defn}

به شرط آنکه $\delta_1$ و $\delta_2$ !!{formula}های فرضی متفاوتی با
برچسب یکسان نداشته باشند و هیچ فرض باز~$\delta_2$ شامل متغیر ویژه‌ای
از~$\delta_1$ نباشد، حاصل
$\Subst{\delta_1}{\delta_2}{!B^x}$ یک !!{proof} درست است و فرض‌های باز
آن، فرض‌های باز~$\delta_1$ همراه با فرض‌های باز~$\delta_2$ هستند.
هنگام پیوندزدن !!{proof}ها، این شرط‌ها معمولاً برقرارند. شرط نخست هرگاه
$\delta_1$ و $\delta_2$ زیربرهان‌های !!{proof} بزرگ‌تری باشند برقرار
است. اگر شرط دوم برقرار نباشد، می‌توان متغیرهای ویژهٔ~$\delta_1$ را
تغییر نام داد تا برقرار شود.

تعریف متناظری برای \Log{N2c} و~\Log{N2i} برقرار است.

\begin{defn}
اگر $\delta_1$ یک \Log{N2c}-!!{proof} (یا \Log{N2i}-!!{proof}) برای
~$x:!B, \Gamma_1 \Sequent !A$، و $\delta_2$ یک
\Log{N2c}-!!{proof} (یا \Log{N2i}-!!{proof}) برای
~$\Gamma_2 \Sequent !B$ باشد، آنگاه
$\Subst{\delta_1}{\delta_2}{x:!B}$ حاصل جایگزینی هر اصل
$x:!B \Sequent !B$ در~$\delta_1$ با~$\delta_2$ و افزودن~$\Gamma_2$
به بافت هر دنبالهٔ زیر آن اصل‌هاست.
\end{defn}

\begin{prop}
اگر $\delta_1 \Proves x:!B, \Gamma_1 \Sequent !A$ و
$\delta_2 \Proves \Gamma_2 \Sequent !B$، آنگاه
$\Subst{\delta_1}{\delta_2}{x:!B}
 \Proves \Gamma_1, \Gamma_2 \Sequent !A$، مشروط بر اینکه هیچ متغیر
ویژهٔ~$\delta_1$ در~$\Gamma_2$ ظاهر نشود.
\end{prop}

\begin{proof}
با استقرا بر $\pheight{\delta_1}$.
\end{proof}

\end{document}
